Tous ces résultats sont unis parce qu’ils transforment la valeur terminale en l’expression quantitative d’une structure préalablement construite. L’exactitude numérique accompagne une conception différente de ce qu’est une constante fondamentale. La connexion nonadique n’engendre pas d’abord trois constantes pour en ajouter ensuite une quatrième : le prolongement \(w_6\to w_{12}\to w_{18}\to w_{24}\to w_{30}\to R_{36}\to G_9\) stabilise un état coinductif commun à partir duquel \(\pi\), \(e\), \(\varphi\) et \(\alpha\) sont publiés de manière corrélée. Leurs lecteurs distinguent clôture, propagation, auto-échelle et couplage ; la coordonnée \(\alpha\) conserve, au sein de cette génération conjointe, sa réalisation signée et dodécaphasique et ses deux voies internes de dérivation.\par

Dans la description habituelle, une constante est souvent présentée comme un nombre placé devant une équation : une quantité dont la théorie a besoin, mais dont la valeur doit être reçue de l’extérieur. Dans l’holographie modulaire triadique, c’est l’inverse qui se produit. La constante apparaît à la fin, lorsqu’un même état a été parcouru par plusieurs lecteurs et que l’on vérifie quelle relation leur a tous survécu.\par

C’est pourquoi les constantes du document sont des grandeurs dotées d’une généalogie : des restes invariants d’une transformation. Chacune dit ce qui a été conservé lorsque quelque chose a changé :\par

\begin{itemize}[leftmargin=*,itemsep=0.35em]
\item un feuillet a été inversé ;
\item un cycle est revenu à sa phase initiale ;
\item une mémoire a traversé la couture ;
\item une description intérieure a été publiée au bord ;
\item une grandeur adimensionnelle a reçu une réalisation physique ;
\item un système a changé d’échelle en conservant intégralement sa généalogie.
\end{itemize}

La mécanique dimensionnelle ajoute ensuite les droites de grandeur et les unités. Mais la structure qui sélectionne la constante est antérieure. La valeur réelle apparaît à la fin parce que \(\mathbb R\) fonctionne ici comme évaluation terminale : elle est l’ombre quantitative d’une construction qui était déjà complète avant de devenir un nombre décimal.\par

Dans cette perspective, les neuf développements considérés forment une seule narration sur l’orientation, la mémoire, la géométrie, l’information et l’échelle.\par
